\chapter{Background: Diffusion as Iterative Denoising}
\label{ch:background-diffusion}

Diffusion models are generative models built around a simple idea: corrupt data
gradually, then learn to reverse the corruption. This chapter reviews that idea
at the level needed for the thesis. The goal is not to develop the full
continuous-time theory of score-based diffusion, but to isolate the concepts
that survive the move to text: a state space, a noising process, a learned
reverse process, a denoising horizon, and an inference-time budget.

\section{Generative Modelling Setup}
\label{sec:ch2-generative-setup}

Let $\pdata$ denote the unknown data distribution on a space $\mathcal{X}$.
A generative model defines a distribution $p_\theta$ intended to approximate
$\pdata$. Sampling from $p_\theta$ should produce new objects that resemble the
training data. In language modelling, the objects are token sequences; in
image modelling, they are arrays of pixel values or latent vectors.

Diffusion models introduce an auxiliary family of noisy distributions
$(p_t)_{t\in[0,1]}$. The endpoint $p_0$ is the data distribution and $p_1$ is a
simple noise distribution. Generation starts from noise and follows a learned
reverse path back toward data. In discrete implementations the interval
$[0,1]$ is replaced by a finite horizon $T$ and a sequence of states
$x_T,x_{T-1},\ldots,x_0$.

The modelling problem has two parts. The forward process $q$ is chosen by the
designer and can be simulated exactly. The reverse process is learned. At
sampling time, the model uses the learned reverse process to denoise a noisy
state step by step.

\section{Forward Noising and Reverse Denoising}
\label{sec:ch2-forward-reverse}

In continuous diffusion models, the forward process is often Gaussian. A clean
data point $x_0$ is progressively mixed with noise:
\[
  x_t = \sqrt{\bar{\alpha}_t}\,x_0
        + \sqrt{1-\bar{\alpha}_t}\,\varepsilon,
  \qquad \varepsilon \sim \mathcal{N}(0,I).
\]
When $t$ is large, the signal term is small and $x_t$ is nearly Gaussian noise.
The DDPM formulation of \textcite{ho2020ddpm} uses this idea in discrete time.
Score-based diffusion models express the same principle through stochastic
differential equations and learned score functions \parencite{song2021score}.

The reverse step asks for a distribution of less-noisy states given the current
noisy state. In general this reverse conditional is not available exactly, so a
neural network is trained to approximate the relevant denoising quantity. One
may train the model to predict the added noise, the clean data, or the score
$\nabla_x \log p_t(x)$, depending on parameterisation
\parencite{song2019ncsn,vincent2011connection}. All three views support the
same operational picture: generation is iterative denoising.

\section{Markov-Chain Viewpoint}
\label{sec:ch2-markov-view}

The sampling procedure can be viewed as an inhomogeneous Markov chain. A state
at time $t$ is transformed into a state at time $t-1$ using a transition kernel
that depends on the current noise level:
\[
  x_t \longrightarrow x_{t-1} \longrightarrow \cdots \longrightarrow x_0 .
\]
The transition kernel is approximate because it uses a learned denoiser and a
finite discretisation. The number of transitions is an inference-time design
choice. More transitions usually means more computation and a closer
approximation to the ideal reverse process; fewer transitions reduce compute
but increase discretisation or approximation error.

This Markov-chain viewpoint is important for the thesis because correctors are
additional transition kernels inserted into an existing trajectory. In
continuous predictor-corrector samplers, the predictor is a discretised reverse
SDE or ODE step and the corrector is often a Langevin step that refines the
current state \parencite{song2021score}. In discrete diffusion, the analogue is
not a gradient update but a local Markov update on tokens.

\section{Continuous and Discrete Diffusion}
\label{sec:ch2-continuous-discrete}

The continuous setting is mathematically convenient because the state space is
Euclidean. Small perturbations are meaningful, gradients exist, and Gaussian
transition kernels are tractable. Text is different. A token sequence lies in
the finite product space $\mathcal{V}^L$, where $\mathcal{V}$ is the
vocabulary and $L$ is the sequence length. There is no infinitesimal move from
one word to another.

Discrete diffusion models therefore replace Gaussian noising with categorical
corruption. Some models substitute tokens according to a rate matrix; masked
diffusion models use a special absorbing mask token $\Mask$. A clean token is
either preserved or replaced by $\Mask$, and the reverse model predicts the
missing content. This connects diffusion language models to masked language
modelling, but with an explicit generative trajectory rather than a single
training-time corruption.

The move to discrete states changes the inference problem. In continuous
models, a denoising network can define a vector field. In masked language
models, the denoiser returns categorical probabilities over the vocabulary at
masked or revisable positions. Sampling requires committing to tokens. Once a
token is committed, later mechanisms may or may not be allowed to revise it,
depending on the sampler.

\section{Inference-Time Budgets}
\label{sec:ch2-budgets}

Diffusion generation is expensive because it uses repeated model evaluations.
The usual accounting unit is the number of function evaluations, or NFE. A
shorter denoising trajectory uses fewer NFEs but may produce lower-quality
samples. Deterministic samplers such as DDIM were introduced partly to improve
this tradeoff in continuous diffusion \parencite{song2020ddim}.

The thesis uses a more specific budget. The base predictor trajectory is fixed,
so its cost is treated as given. The variable cost is the number of corrector
placements. If a corrector placement costs $c_{\mathrm{corr}}$ extra model
forward passes, then a placement budget $B$ corresponds to extra compute
$B_{\mathrm{NFE}} = c_{\mathrm{corr}}B$. The mathematical problem is therefore
not generic NFE allocation; it is placement of a fixed number of correction
calls along an existing denoising trajectory.

\section{Quality Functionals and Paired Evaluation}
\label{sec:ch2-quality}

To compare two generation procedures one needs a quality functional. In this
thesis a quality functional is any fixed map $F$ from a completed sequence to a
real number, with larger values treated as better. Examples include negative
language-model loss under a reference model, learned reward scores, or other
automatic metrics. The theory in Chapter~\ref{ch:contribution} is agnostic to
the choice of $F$; the empirical conclusions are always conditional on the
chosen $F$.

Because generation is random, schedule comparisons should be paired whenever
possible. A paired comparison holds the seed and baseline trajectory fixed and
changes only the intervention being tested. This common-random-number
principle reduces variance and is central to the schedule-gain definition used
later.

\section{Summary}
\label{sec:ch2-summary}

Diffusion models generate by iterative denoising. The details differ between
continuous and discrete state spaces, but the common structure is a noisy
starting state, a learned reverse process, a finite denoising horizon, and an
inference-time budget. For text, the natural diffusion state is a partially
masked token sequence. The next chapter develops this discrete masked setting
in detail.
